\clearpage

\addcontentsline{toc}{chapter}{Abstract}
\begin{center}
	\textbf{\large Abstract}
\end{center}

Named Entity Recognition (NER) in new or specialized domains often requires manually defining entity types and annotating data, a process that is costly and demands domain expertise. 
This limitation motivates \textbf{open-world NER}, where the goal is not only to detect entity mentions from unlabeled text but also to \textbf{discover and name latent entity types} automatically. 
In this thesis, we propose \textbf{SAGE}, a semantic-augmented, graph-based framework for open-world entity typing and ontology induction.

SAGE operates without gold entity labels in the target corpus. It first obtains high-recall entity mentions by aggregating predictions from multiple distilled NER teachers, 
such as GLiNER and UniversalNER, using a small, domain-agnostic label bank solely for span proposal. Each mention is then represented by a \textbf{type signature} that combines contextual semantic 
embeddings, teacher-generated type phrases, and corpus-grounded type evidence mined from definitional and Hearst-style patterns. Using these signatures, SAGE constructs a similarity graph over mentions and applies constraint-aware community detection to induce an initial set of entity type clusters without predefining the number of types.

To address common clustering failures in open-world settings, we introduce a \textbf{Critic refinement module} that evaluates cluster quality through learnability. 
A lightweight classifier is trained to predict cluster assignments, and its confusion patterns and prediction entropy are used to identify over-merged and internally inconsistent clusters, 
which are iteratively merged or split in a label-free manner. Finally, SAGE generates human-readable entity type names by mining candidate phrases from clusters and selecting concise, 
specific labels using an open-source language model.

Experiments on multiple public benchmarks across general, noisy, and biomedical domains demonstrate that SAGE improves type discovery quality and produces more accurate and 
stable entity class names compared to clustering-only and teacher-only baselines. These results show that SAGE provides a practical and effective solution 
for automatically constructing entity type inventories when deploying NER systems in new domains.


\vspace*{1cm}

\noindent\textbf{Keywords:} Open-world Named Entity Recognition; Entity Type Discovery; Ontology Induction; Entity Class Name Generation; Graph-based Clustering; Teacher-Augmented Learning; GLiNER; Universal NER; Self-supervised Refinement; Learnability-based Clustering; Domain Adaptation

\begin{comment}


\end{comment}


